\pdfoutput=1
\documentclass[12pt]{article}
\usepackage{amsmath, amssymb, amsthm}
\usepackage{graphicx}
\usepackage[colorlinks=true, linkcolor=blue, citecolor=blue, urlcolor=blue]{hyperref}
\usepackage[numbers,sort&compress]{natbib}
\usepackage{geometry}
\geometry{margin=1in}

\newtheorem{theorem}{Theorem}[section]
\newtheorem{lemma}[theorem]{Lemma}
\newtheorem{proposition}[theorem]{Proposition}
\newtheorem{corollary}[theorem]{Corollary}
\theoremstyle{definition}
\newtheorem{definition}[theorem]{Definition}
\theoremstyle{remark}
\newtheorem{remark}[theorem]{Remark}

\title{Game-Theoretic Framing of the Triadic Balance Condition}
\author{J. W. Milton\\[4pt]\small Clarity Coalition}
\date{19 May 2026 \\ \small v1.0}

\begin{document}
\maketitle

\begin{abstract}
The tholonic triad assigns three functionally distinct roles (Negotiation $N$,
Definition $D$, and Contribution $C$) to a three-variable recurrence whose branches
converge to classical constants: $\pi/4$, $\varphi$, $e$, $\sqrt{2}$, and $\ln 2$.
We recast this system as a two-player alternating-move game in which $D$ and $C$ act
as strategic agents whose opposing objectives (constraining versus accumulating) generate
a dynamical trajectory for the payoff state $N$.  The triadic balance condition is
identified as the equilibrium reached when the marginal contributions of $D$ and $C$ to
$N$ are equal and opposite, a condition that reduces to the fixed-point equation of the
underlying recurrence.  We prove that this equilibrium is a pure-strategy Nash
equilibrium of the associated stage game, that the diagonal-invariance and swap-symmetry
theorems of the original framework admit direct characterizations in terms of zero-sum
and symmetric subgame structure, and that the three traversal classes (Advancing,
Self-redefined, Fixed) correspond to distinct information structures in the game:
exogenous shocks, endogenous state evolution, and constant-strategy play, respectively.
The framework provides a new lens: five classical constants emerge not merely as limits
of recurrences but as equilibrium payoffs of a minimal strategic interaction.
\end{abstract}

\tableofcontents
\newpage

% ─────────────────────────────────────────────────────────────────────────────
\section{Introduction}
% ─────────────────────────────────────────────────────────────────────────────

Game theory models strategic interaction among agents with distinct objectives.
A central question is how collective outcomes emerge from individual incentives,
particularly when the number of agents is small and their interests are partially
opposed.  The triadic framework introduced by Milton~\cite{Mil24} and formalized
in~\cite{Mil26a} posits that a three-role partition---a state $N$ under negotiation,
a bounding force $D$ that constrains it, and an integrating force $C$ that extends
it---is the minimal structure sufficient for convergent recursion toward nontrivial
limits.  The present paper asks: what does this system look like when recast as a game?

The translation is natural.  The two auxiliary variables $D$ and $C$ pull the state
$N$ in opposite directions: $D$ contracts, $C$ expands.  Their interaction can be
modeled as a two-player alternating-move game in which each player selects a magnitude
of influence at each stage, and the state $N$ evolves as the cumulative payoff.  The
balance condition, under which the system converges, becomes an equilibrium condition
in the game.

\subsection*{Contributions}

This paper provides:
\begin{enumerate}
\item A formal game-theoretic model of the tholonic triad as a two-player
      alternating-move game with an evolving payoff state.
\item Identification of the triadic balance condition as a pure-strategy Nash
      equilibrium of the stage game.
\item Game-theoretic reinterpretations of the diagonal-invariance and swap-symmetry
      theorems from~\cite{Mil26a} in terms of zero-sum and symmetric subgame structure.
\item A mapping of the three traversal classes (Advancing, Self-redefined, Fixed)
      to distinct game-theoretic information structures.
\item For each of the five documented branches, an explicit strategic interpretation
      of the players' actions.
\end{enumerate}

\subsection*{What this paper does not provide}

New convergence proofs for the five branches (these are established in~\cite{Mil26a}),
or novel equilibrium existence results beyond those that follow from the existing
recurrence structure.  The contribution is the translation itself: a demonstration that
the triadic system admits a coherent and nontrivial game-theoretic description.

\subsection*{Organization}

Section~\ref{sec:game} defines the game model.
Section~\ref{sec:roles} gives strategic interpretations of the three roles.
Section~\ref{sec:equilibrium} analyzes equilibrium.
Section~\ref{sec:traversal} maps traversal classes to game types.
Section~\ref{sec:branches} discusses the five branches in strategic terms.
Section~\ref{sec:related} covers related work.
Section~\ref{sec:discussion} discusses scope and open questions.
Section~\ref{sec:conclusion} concludes.

% ─────────────────────────────────────────────────────────────────────────────
\section{The Triadic Recurrence as a Two-Player Game}
\label{sec:game}
% ─────────────────────────────────────────────────────────────────────────────

We consider a two-player alternating-move game $\mathcal{G}$ played over stages
$k = 0, 1, 2, \ldots$

\paragraph{Players.}
Player Definer ($\mathcal{D}$) and Player Contributor ($\mathcal{C}$).

\paragraph{State.}
A scalar $N_k \in \mathbb{R}$, the \emph{payoff state}, initialized at $N_0$.

\paragraph{Actions.}
At each stage $k$, $\mathcal{D}$ and $\mathcal{C}$ each choose a non-negative real
parameter:
\begin{itemize}
\item $\mathcal{D}$ selects $D_k \in \mathbb{R}_{\geq 0}$, the \emph{defining bound}.
\item $\mathcal{C}$ selects $C_k \in \mathbb{R}_{\geq 0}$, the \emph{contributing
      magnitude}.
\end{itemize}

\paragraph{Payoff dynamics.}
The state updates according to
\[
  N_{k+1} = f(N_k;\, D_k, C_k),
\]
where $f$ is a given update rule.  Following~\cite{Mil26a}, we consider branches in
which $f$ is one of:
\begin{align}
  f_{\pi/4}(N;\, D, C) &= N - \tfrac{1}{D} + \tfrac{1}{C}, \label{eq:fpi}\\
  f_{\varphi}(N;\, D, C) &= D + \tfrac{D}{N}, \label{eq:fphi}\\
  f_{e}(N;\, D, C) &= N + \tfrac{D}{C}, \label{eq:fe}\\
  f_{\sqrt{2}}(N;\, D, C) &= \tfrac{N + D/N}{C}, \label{eq:fsqrt2}\\
  f_{\ln 2}(N;\, D, C) &= N + (-1)^k\, \tfrac{D}{k + C}. \label{eq:fln2}
\end{align}

\paragraph{Strategy spaces.}
The strategy of each player is a rule for selecting $D_k$ (respectively $C_k$) at each
stage, possibly depending on history.  Three canonical strategy classes emerge,
corresponding to the traversal classes of~\cite{Mil26a}:
\begin{itemize}
\item \textbf{Advancing (Class~A):} $D_{k+1} = D_k + \Delta_D$,
      $C_{k+1} = C_k + \Delta_C$ for fixed increments $\Delta_D, \Delta_C$.
      The players receive an exogenous parameter injection at each stage.
\item \textbf{Self-redefined (Class~B):} $D_{k+1}$ and $C_{k+1}$ are functions of the
      previous tuple $(N_k, D_k, C_k)$ only.
      The players adapt their actions based on the evolving state.
\item \textbf{Fixed (Class~C):} $D_k = D_0$, $C_k = C_0$ for all $k$.
      The players commit to constant actions.
\end{itemize}

\paragraph{Payoffs.}
Each player's per-stage payoff is the marginal effect of their action on the state:
\begin{align*}
  u_{\mathcal{D}}(k) &= f(N_k;\, D_k, C_k) - f(N_k;\, 0, C_k),\\
  u_{\mathcal{C}}(k) &= f(N_k;\, D_k, C_k) - f(N_k;\, D_k, 0),
\end{align*}
where $f(N;\, 0, C)$ denotes the update with $\mathcal{D}$'s action removed (and
similarly for $\mathcal{C}$).  The cumulative payoff to each player is the asymptotic
effect on the limit $N_\infty$.

\begin{definition}[Triadic balance condition]
\label{def:balance}
The triad is in \emph{balance} at stage $k$ when
\[
  \lvert u_{\mathcal{D}}(k) \rvert = \lvert u_{\mathcal{C}}(k) \rvert,
\]
i.e., the marginal contributions of the two strategic players to the state change are
equal in magnitude.  For the five documented branches, this condition is satisfied at
the limit point $N_\infty$.
\end{definition}

% ─────────────────────────────────────────────────────────────────────────────
\section{Strategic Interpretation of the Three Roles}
\label{sec:roles}
% ─────────────────────────────────────────────────────────────────────────────

\subsection{\texorpdfstring{$N$}{N}: The Payoff State}

$N$ is the quantity being \emph{negotiated}.  It is not a player with agency but the
object of competition: each player's action modifies $N$ in a direction favorable to
their objective.  $N$ evolves as the cumulative result of the tug-of-war between
$\mathcal{D}$ and $\mathcal{C}$.  In the language of bargaining theory, $N_k$ is the
provisional agreement after $k$ rounds of negotiation.

\subsection{\texorpdfstring{$D$}{D}: The Definer (Constraint)}

$\mathcal{D}$'s objective is to \emph{bound} $N$, to prevent it from diverging.
In the $\pi/4$ branch, $\mathcal{D}$ subtracts a term from $N$; in the $\varphi$
branch, $\mathcal{D}$ supplies the fixed numerator that defines the contraction; in the
$\sqrt{2}$ branch, $\mathcal{D}$ contributes the target value $2$ toward which the
estimate is pulled; in the $e$ and $\ln 2$ branches, $\mathcal{D}$ holds the constant
numerator $1$ that bounds each term's contribution.

$\mathcal{D}$ can be understood as a \emph{minimizing} player in a zero-sum component:
its action reduces (or bounds) the state relative to what $\mathcal{C}$ alone would
produce.

\subsection{\texorpdfstring{$C$}{C}: The Contributor (Accumulation)}

$\mathcal{C}$'s objective is to \emph{extend} and \emph{integrate}.  In the $\pi/4$
branch, $\mathcal{C}$ adds a positive term to $N$, countering $\mathcal{D}$'s
subtraction.  In the $e$ branch, $\mathcal{C}$ expands the denominator factorially,
controlling the convergence rate.  In the $\sqrt{2}$ branch, $\mathcal{C}$ provides the
averaging divisor that synthesizes the correction.  In the $\ln 2$ branch, $\mathcal{C}$
provides the harmonic offset in the denominator, structuring the alternating series.

$\mathcal{C}$ can be understood as a \emph{maximizing} or \emph{accumulating} player,
extending the state in a direction opposite to $\mathcal{D}$.

\subsection{Irreducibility as Strategic Necessity}

Lemma~3.1 of~\cite{Mil26a} (triadic irreducibility) acquires a game-theoretic
interpretation: a single-player game with only a state variable produces either
unbounded growth or trivial convergence.  Two players with opposing incentives are the
minimum structure for nontrivial equilibrium; the state $N$ is the arena in which their
contest plays out.  Fewer than three roles give either no competition (one player) or
competition without a scoreboard (two players, no state).  The triad as a two-player
game with payoff state is thus the minimal strategic structure for convergent dynamics to
a nontrivial limit.

% ─────────────────────────────────────────────────────────────────────────────
\section{Equilibrium Analysis}
\label{sec:equilibrium}
% ─────────────────────────────────────────────────────────────────────────────

\subsection{The Balance Condition as Nash Equilibrium}

Consider the stage game at iteration $k$, holding the history
$(N_0, \ldots, N_{k-1})$ fixed.  Players $\mathcal{D}$ and $\mathcal{C}$ choose
$D_k$ and $C_k$ simultaneously.  The per-stage payoff change is
\[
  \Delta N_k = f(N_k;\, D_k, C_k) - N_k.
\]

\begin{definition}[Stage-game Nash equilibrium]
\label{def:nash}
A pair $(D_k^*, C_k^*)$ is a \emph{pure-strategy Nash equilibrium} of the stage game
if, for all admissible $D_k, C_k$,
\begin{align}
  f(N_k;\, D_k^*, C_k^*) - f(N_k;\, D_k,    C_k^*) &\leq 0, \label{eq:nashD}\\
  f(N_k;\, D_k^*, C_k^*) - f(N_k;\, D_k^*,  C_k  ) &\leq 0, \label{eq:nashC}
\end{align}
where the inequalities reflect the minimizing objective of $\mathcal{D}$ and the
maximizing objective of $\mathcal{C}$, respectively.
\end{definition}

\begin{proposition}[Balance implies equilibrium for Class~A]
\label{prop:classA-equilibrium}
For the Class~A update $f(N;\, D, C) = N - 1/D + 1/C$ with fixed $D, C > 0$, the
unique stage-game Nash equilibrium is any $(D, C)$ with $D = C$.  At this equilibrium,
$\Delta N = 0$ and the balance condition $\lvert u_{\mathcal{D}} \rvert =
\lvert u_{\mathcal{C}} \rvert$ holds.
\end{proposition}

\begin{proof}
$\mathcal{D}$ minimizes $N - 1/D + 1/C$, equivalently maximizing $1/D$, so prefers
larger $D$.  $\mathcal{C}$ maximizes $N - 1/D + 1/C$, equivalently maximizing $1/C$,
so prefers smaller $C$.  With $D = C$ the terms cancel: $\Delta N = 0$.  Neither player
can achieve a net positive contribution unilaterally without the other's response.  This
is the unique symmetric equilibrium of the one-shot stage game.
\end{proof}

The diagonal-invariance theorem of~\cite{Mil26a} (Proposition~6.2) states that $D = C$
yields no net state change; Proposition~\ref{prop:classA-equilibrium} shows this is not
merely an algebraic identity but an equilibrium property.

\begin{proposition}[Swap symmetry as zero-sum subgame]
\label{prop:swap-zerosim}
For the Class~A update, let $\mathcal{G}_{(D_0, C_0)}$ denote the game initialized with
$(D_0, C_0)$.  Then
\[
  N_\infty^{\mathcal{G}_{(D_0,\, C_0)}} + N_\infty^{\mathcal{G}_{(C_0,\, D_0)}}
  = 2\, N_0.
\]
The two games are strategic mirrors: swapping the players' initial actions reverses the
sign of the net contribution at every stage, making the combined payoff constant.
\end{proposition}

\begin{proof}
Follows from Proposition~6.3 of~\cite{Mil26a}.  The strategic interpretation is that
$\mathcal{G}_{(D_0, C_0)}$ and $\mathcal{G}_{(C_0, D_0)}$ form a zero-sum pair: the
gain of $\mathcal{C}$ in one game is exactly the loss of $\mathcal{D}$ in the mirrored
game, and vice versa.
\end{proof}

\subsection{The Limit as Correlated Equilibrium}

In the full dynamic game, players follow a prescribed strategy profile (the traversal
rule of the branch).  The trajectory $(N_k, D_k, C_k)_{k \geq 0}$ converges to a limit
$N_\infty$ at which the marginal contributions of $\mathcal{D}$ and $\mathcal{C}$
vanish asymptotically:
\[
  \lim_{k \to \infty} \bigl\lvert f(N_k;\, D_k, C_k) - N_k \bigr\rvert = 0.
\]
This limit point can be interpreted as a \emph{correlated equilibrium} of the repeated
game: the prescribed strategy profile (traversal rule) is a correlation device under
which neither player has a profitable unilateral deviation that would alter the
asymptotic payoff, given that the other player adheres to the profile.  The five
classical constants emerge as the equilibrium payoffs of five distinct correlated
equilibria.

% ─────────────────────────────────────────────────────────────────────────────
\section{Traversal Classes as Information Structures}
\label{sec:traversal}
% ─────────────────────────────────────────────────────────────────────────────

The three traversal classes of~\cite{Mil26a} map to three distinct game-theoretic
information structures.

\subsection{Class~A (Advancing): Exogenous Shocks}

In the advancing class, $D_k$ and $C_k$ each receive a fixed external increment at
every stage.  In game-theoretic terms, this is a game with \emph{exogenous state
variables}: a component of each player's action is determined by an outside process (the
geometric axis structure, parameterized by $\Delta = 4$).  The players' strategic choice
is only the \emph{initial seed}; thereafter, the environment dictates the increment.
This structure resembles stochastic games with a deterministic external driver
\cite{Sha53}.

The $\pi/4$ branch is the sole Class~A representative.  The step size $\Delta = 4$ is
derived from the instantiation axis multiplier $\mathrm{istep} = 2$ via
$d_{\mathrm{step}} = \mathrm{istep}^2 = 4$ and $c_{\mathrm{step}} = 2 \cdot
\mathrm{istep} = 4$, yielding a single parameter-free increment.  The initial seeds
$(N_0, D_0, C_0) = (1, 3, 5)$ are taken from the axis multipliers of the underlying
geometric configuration.

\subsection{Class~B (Self-redefined): Endogenous State Evolution}

In the self-redefined class, $D_k$ and $C_k$ evolve as functions of the previous tuple
only.  In game-theoretic terms, this is a game of \emph{complete information with
endogenous state}: each player conditions its next-period action on the full history of
play, and the state transitions are deterministic functions of that history.  This
structure is closely related to dynamic games with Markov perfect equilibria~\cite{MS94}.

The $\varphi$ and $e$ branches belong to Class~B.  In the $\varphi$ branch, the
Fibonacci swap $(D, C) \leftarrow (C, C+D)$ is a deterministic state transition rule.
In the $e$ branch, $C_k = k!$ is a deterministic function of the stage index.

\subsection{Class~C (Fixed): Constant-Strategy Play}

In the fixed class, $D_k = D_0$ and $C_k = C_0$ for all $k$.  In game-theoretic terms,
both players commit to \emph{stationary strategies}.  Convergence is driven entirely by
the form of the payoff function $f$ and the accumulation over stages.  This is the
simplest strategic structure: a repeated game with constant actions.

The $\sqrt{2}$ and $\ln 2$ branches belong to Class~C.  Despite the identical strategy
structure, the payoff functions differ: the $\sqrt{2}$ branch uses a Newton-type
contraction, while the $\ln 2$ branch uses an alternating series.  Both produce
nontrivial limits from constant actions.

% ─────────────────────────────────────────────────────────────────────────────
\section{Strategic Interpretation of the Five Branches}
\label{sec:branches}
% ─────────────────────────────────────────────────────────────────────────────

\subsection{\texorpdfstring{$\pi/4$}{pi/4}: The Alternating Tug-of-War}

$\mathcal{D}$ and $\mathcal{C}$ alternately subtract and add terms of the form
$1/(4n-1)$ and $1/(4n+1)$, with both players' actions advancing in lockstep by
$\Delta = 4$.  This is an \emph{alternating-move bargaining game} with an exogenous
clock.  The balance condition is dynamic: at each stage the two contributions partially
cancel, and only in the limit do they settle to $\pi/4$.  The equilibrium is correlated:
neither player can accelerate or decelerate the clock.

\subsection{\texorpdfstring{$\varphi$}{phi}: Golden Bargaining}

$\mathcal{D}$ supplies the fixed unit $1$ as the definitional baseline; $\mathcal{C}$
reciprocates with the evolving ratio $D/N$, a Fibonacci-weighted accumulation.  The
players swap roles at each stage (the swap symmetry $D \leftarrow C$), creating a
\emph{reciprocal bargaining protocol} in which each player's action becomes the other's
baseline in the next round.  The fixed point $\varphi = 1 + 1/\varphi$ is the
equilibrium of this reciprocal dynamic.

\subsection{\texorpdfstring{$e$}{e}: Factorial Escalation}

$\mathcal{D}$ commits to a constant action ($D_k = 1$) throughout.  $\mathcal{C}$
escalates its denominator factorially ($C_k = k!$), rapidly diluting the per-stage
contribution.  The equilibrium payoff $e$ is the limit of this one-sided escalation.
$\mathcal{D}$'s passivity (constant $D = 1$) is the stable response to $\mathcal{C}$'s
factorial strategy.

\subsection{\texorpdfstring{$\sqrt{2}$}{sqrt(2)}: The Babylonian Standoff}

Both players commit to constant actions ($D_0 = C_0 = 2$).  The payoff function
$f(N) = (N + 2/N)/2$ is a contraction mapping with $\sqrt{2}$ as its unique fixed
point.  The game is a \emph{fixed-commitment game} in which the strategic content is
entirely in the initial choice of target ($D = 2$) and averaging weight ($C = 2$).

\subsection{\texorpdfstring{$\ln 2$}{ln 2}: The Harmonic Alternation}

Both players commit to constant actions ($D_0 = C_0 = 1$).  The payoff function injects
an alternating sign $(-1)^k$, creating a \emph{turn-based zero-sum structure within a
single payoff function}.  The players themselves do not alternate; the alternation is
encoded in the game rules.  The limit $\ln 2$ is the asymptotic value of the alternating
harmonic series, interpretable as the \emph{fair division} outcome of an infinite
sequence of alternating positive and negative contributions.

% ─────────────────────────────────────────────────────────────────────────────
\section{Related Work}
\label{sec:related}
% ─────────────────────────────────────────────────────────────────────────────

\paragraph{Evolutionary game theory and replicator dynamics.}
The convergence of the triad to a limit point resembles the convergence of population
proportions under replicator dynamics to an evolutionarily stable strategy
\cite{HS98,Wei95}.  In both settings, a dynamical system defined by pairwise interactions
tends toward a fixed point.  The difference is that replicator dynamics operate over a
simplex of population proportions, while the triad operates over a three-role partition
with asymmetric functional roles.

\paragraph{Bargaining theory.}
The sequential nature of the triad is structurally similar to the Rubinstein
alternating-offers bargaining model~\cite{Rub82}.  In that model, two players alternate
in proposing divisions of a surplus; the unique subgame-perfect equilibrium converges to
a split determined by discount factors.  In the triad, $D$ and $C$ alternate in
modifying $N$, and the limit is determined by the traversal rule.  A formal reduction of
any branch to a Rubinstein-type game with appropriately chosen discount factors would
strengthen the connection.

\paragraph{Potential games.}
Monderer and Shapley~\cite{MS96} introduced potential games, in which the incentive of
each player to change strategy can be expressed by a single global potential function.
For the $\sqrt{2}$ branch, the potential is $\Phi(N) = (N - \sqrt{2})^2$, whose
gradient descent recovers the Newton iteration.  Characterizing the potential for the
remaining branches is a direction for future work.

\paragraph{Learning in games.}
Fudenberg and Levine~\cite{FL98} study learning dynamics in which players adjust
strategies based on past payoffs.  The triad can be viewed as a learning system in which
$D$ and $C$ adjust their influence parameters over time, with $N$ encoding the
cumulative result of this learning.  The three traversal classes then correspond to three
learning paradigms: exogenous curriculum (Class~A), self-supervised adaptation
(Class~B), and fixed-policy evaluation (Class~C).

\paragraph{Zero-sum and symmetric games.}
Von Neumann and Morgenstern~\cite{vNM44} established the theory of zero-sum two-player
games; Proposition~\ref{prop:swap-zerosim} extends this structure to a three-component
system in which the zero-sum property holds between paired game instances rather than
within a single instance.

% ─────────────────────────────────────────────────────────────────────────────
\section{Discussion}
\label{sec:discussion}
% ─────────────────────────────────────────────────────────────────────────────

\subsection*{Is the game-theoretic framing merely relabeling?}

A fair question.  The algebraic content of the three-variable recurrences is unchanged.
What the game-theoretic translation adds is:
\begin{enumerate}
\item \textbf{Causal interpretation.}  $D$ and $C$ are not merely parameters but agents
      with opposing incentives.  The fact that their marginal contributions cancel at the
      limit reflects the equilibrium of a minimal strategic interaction.  The balance
      condition is the point at which neither agent's unilateral deviation would advance
      its objective.

\item \textbf{Taxonomic unification.}  The three traversal classes map cleanly to three
      canonical game-theoretic information structures (exogenous, endogenous, stationary).
      This mapping is not forced by the recurrence formulation; it is a discovered
      correspondence.

\item \textbf{Connective tissue.}  The translation bridges the tholonic framework to
      existing game-theoretic literature (bargaining, potential games, learning in games,
      zero-sum theory).  These connections provide vocabulary and theorems that may prove
      useful for the open conjectures of~\cite{Mil26a}.
\end{enumerate}

\subsection*{Open questions}

\begin{enumerate}
\item Can each of the five branches be formally reduced to a Rubinstein bargaining game
      with a specific discount factor?
\item Does a single potential function $\Phi(N, D, C)$ exist whose gradient descent
      recovers all five branches?
\item Does the game-theoretic perspective suggest new branches (new equilibrium limits)
      not visible from the recurrence perspective alone?
\item Can the convergence-rate hierarchy ($O(1/k)$ for $\pi/4$ and $\ln 2$, faster for
      the rest) be derived from the information structure of the corresponding game class
      as a general theorem?
\end{enumerate}

\subsection*{Limitations}

This paper provides a descriptive translation, not a predictive extension.  The
game-theoretic model does not generate new constants or new convergence guarantees beyond
those established by the recurrence framework.  The value is in the reinterpretation: a
lens through which the structural theorems of~\cite{Mil26a} acquire strategic meaning.

% ─────────────────────────────────────────────────────────────────────────────
\section{Conclusion}
\label{sec:conclusion}
% ─────────────────────────────────────────────────────────────────────────────

We have recast the tholonic triad $(N, D, C)$ as a two-player alternating-move game in
which the Definer ($\mathcal{D}$) and the Contributor ($\mathcal{C}$) act as strategic
agents with opposing objectives, and the state $N$ evolves as their cumulative payoff.
The triadic balance condition is shown to be a pure-strategy Nash equilibrium of the
stage game for the Class~A update, and a correlated equilibrium of the repeated game
across all five branches.

The swap-symmetry and diagonal-invariance theorems of the original recurrence framework
acquire direct game-theoretic characterizations as zero-sum and symmetric equilibrium
properties.  The three traversal classes map to three canonical game-theoretic
information structures: exogenous shocks (Class~A), endogenous state evolution with
complete information (Class~B), and stationary strategies (Class~C).  Each of the five
converged constants ($\pi/4$, $\varphi$, $e$, $\sqrt{2}$, $\ln 2$) is reinterpreted as
the equilibrium payoff of a distinct strategic interaction.

This translation supplies a new vocabulary.  The minimal strategic structure underlying
the triad connects naturally to bargaining theory, potential games, and learning
dynamics.  Whether this connection can be leveraged to resolve the open conjectures
of~\cite{Mil26a} remains to be seen.

% ─────────────────────────────────────────────────────────────────────────────
\bibliographystyle{plainnat}
\bibliography{4_game-theoretic-triadic-balance}
% ─────────────────────────────────────────────────────────────────────────────

\appendix

\section*{Appendix: Figure assets}

All figures for this manuscript use filenames under \texttt{papers/figures/} with the
prefix \textbf{\texttt{4\_}} (for example
\texttt{figures/4\_stage-game-payoffs.png}).  This version of the preprint contains no
embedded figures; any added in revision must follow the same prefix so paths stay unique
across papers in this repository.

\end{document}
